% page 448 of the facsimile (image p-447 of the scan)
% block 152.4 x 229.0 mm ; left margin 8.0 mm
\pgc{-12.700mm}{0.000mm}{
\l{\cel{3}440.}
\l{}
\l{\sou{piniono}. (\sou{tekn.})\cel{2}Roto dentoza qua ingranas la denti di}
\l{\cel{3}altra roto. - EFS.}
\l{}
\l{\sou{pinselo}. Tufo de pili, ligita an la extremajo di bastono}
\l{\cel{3}mancha, uzata por extensar farbi, gluo, e c . - DFS.}
\l{}
\l{\sou{pinto.}I.Extremajo tale tenuigita ke lu pikas. (Ex.: pinto}
\l{\cel{3}di agulo, di espado, di kompaso). - II. Ulo kelke pinto-}
\l{\cel{3}forma. (ex.: pinto di monto). - EFIS.}
\l{}
\l{\sou{pintad}o. (\sou{zool.}) Ucelo kriema, galinacea, kun kapo nuda,}
\l{\cel{3}kaudo kurta, plumaro griza bluatre, sur qua dissemesas}
\l{\cel{3}makuli blanka. - L.\cel{1}\sou{Numida meleagris.}\cel{1}- EFS.}
\l{}
\l{\sou{piocho.}\cel{2}(\sou{tekn.}) Utensilo di terlaboristo, ek mancho e}
\l{\cel{3}fer-bloko, pikona ye un extremajo, e haua ye la altra. -}
\l{}
\l{}
\l{\sou{piono}. Singla de la disketi di la damo-ludo, e di la fi-}
\l{\cel{3}gurini di la shak-ludo. - FS.}
\l{}
\l{\sou{pioniro}. I. En armeo : laboristo quan onu uzas por voyifar,}
\l{\cel{3}exkavar tranchei, e c. - II. La persono qua kultiveskas}
\l{\cel{3}tereni nekultivita til nun. (\sou{anke metaf.})- DEFIR.}
\l{}
\l{\sou{pipo}. Tubo qua finas per quaza furneleto en qua onu pozas}
\l{\cel{3}tabako quan onu acendas por aspirar lua fumuro +smoke.}
\l{\cel{3}- EFIS.}
\l{}
\l{\sou{pipeto}. (\sou{tekn.)}\cel{3}Tubo vitra, konkavigita ye un extremajo,}
\l{\cel{3}tenuigita til pintesko ye la altra, quan onu povas sto-}
\l{\cel{3}par per un fingro, uzata por transvarsar liquido. - DEFI}
\l{}
\l{\sou{pipiar}. (\sou{netrans.}) (\sou{aludante la ucel-yuni}).Kriar. - DFS.}
\l{}
\l{\sou{pipro.}\cel{2}(\sou{bot.})\cel{2}Mikra bero, qua odoras e saporas forte, e}
\l{\cel{3}qua, sikigite, o muelite, uzesas kom spico, kondimento.}
\l{\cel{3}- EFI.}
\l{}
\l{\sou{pipsar}.\cel{2}(\sou{netrans.})\cel{2}(\sou{aludante la uceli e precipue la hani}).}
\l{\cel{3}Maladesar pro la kresko sur la pinto di la lango, di}
\l{\cel{3}pelikulo squamatra qua impedas drinkar. - DEFIS.}
\l{}
\l{\sou{piquo.}\cel{1}Shaft-armo, di qua la shafto esas garnisita ye ferajo}
\l{\cel{3}plata e pintoza, plu kurta kam lanco. - II. Un ek la du}
\l{\cel{3}figuri nigra, sur ula ludo-karti, di qua la formo esas}
\l{\cel{3}olta di piquo-ferajo. - F.}
\l{}
\l{\sou{piro}.\cel{2}(\sou{bot.}) Frukto di arboro (L.\cel{1}\sou{pirus communis)}\cel{1}ek la fa-}
\l{\cel{3}milio \textquotedbl{}rozacei\textquotedbl{}, oblonga, di qua la volumino deskreskas}
\l{\cel{3}arivante proxim la kaudo. - EFIS.}
}
